% =========================================================================
%  main.tex — Presentazione Tesi di Laurea
%  Base: template beamerthemeUnimore fornito dal correlatore
%  Compilare con: pdflatex main.tex  (eseguire due volte)
% =========================================================================
\documentclass[aspectratio=169]{beamer}
\usetheme{Unimore}

\usepackage[utf8]{inputenc}
\usepackage[italian]{babel}
\usepackage{graphicx}
\usepackage{amsmath}
\usepackage{booktabs}
\usepackage{array}
\usepackage{bm}

% -------------------------------------------------------------------------
% Stili TikZ centralizzati
% Il file tikz_styles.tex si trova nella cartella padre della presentazione.
% Contiene palette, librerie TikZ, nodi e frecce usati in tutte le figure.
% -------------------------------------------------------------------------
% =========================================================================
% CONFIGURAZIONE CENTRALE STILI TIKZ - INGEGNERIA INFORMATICA
% =========================================================================
\usepackage{tikz}
\usetikzlibrary{shapes.geometric, arrows.meta, positioning, calc, fit, backgrounds, shadows}
\usepackage{xcolor}

% -------------------------------------------------------------------------
% 1. PALETTE CROMATICA
% -------------------------------------------------------------------------
\definecolor{mainblue}{HTML}{1F77B4}       % Blu Principale (Input / Primary)
\definecolor{lightblue}{HTML}{E8F1F5}      % Blu Chiaro
\definecolor{accentorange}{HTML}{FF7F0E}  % Arancio (Target / Highlights)
\definecolor{lightorange}{HTML}{FDF0E6}   % Arancio Chiaro
\definecolor{paramgreen}{HTML}{2E7D32}    % Verde (Processing / Fit / Valid)
\definecolor{lightgreen}{HTML}{E8F5E9}    % Verde Chiaro
\definecolor{neutralgrey}{HTML}{4A5568}   % Grigio Scuro (Trasformazioni / Flow)
\definecolor{backgrey}{HTML}{F8FAFC}      % Grigio Chiaro (Output / Files)
\definecolor{alertred}{HTML}{D32F2F}      % Rosso (Anomalie / Attacchi / Errori)
\definecolor{lightred}{HTML}{FFEBEE}      % Rosso Chiaro

% -------------------------------------------------------------------------
% 2. DEFINIZIONE STILI GLOBALI TIKZ (NO RIGHE VUOTE INTERNE)
% -------------------------------------------------------------------------
\tikzset{
    % Font globale e spaziatura interna dinamica
    global_font/.style = {font=\small, align=center, inner sep=6pt},
    % --- NODI RETTANGOLARI STANDARD ---
    node_primary/.style = {
        global_font, rectangle, draw=mainblue, fill=lightblue, thick, rounded corners=4pt,
        drop shadow={shadow xshift=0.12cm, shadow yshift=-0.12cm, opacity=0.12}
    },
    node_accent/.style = {
        global_font, rectangle, draw=accentorange, fill=lightorange, thick, rounded corners=4pt,
        drop shadow={shadow xshift=0.12cm, shadow yshift=-0.12cm, opacity=0.12}
    },
    node_process/.style = {
        global_font, rectangle, draw=paramgreen, fill=lightgreen, thick, rounded corners=4pt,
        drop shadow={shadow xshift=0.12cm, shadow yshift=-0.12cm, opacity=0.12}
    },
    node_transform/.style = {
        global_font, rectangle, draw=neutralgrey, fill=white, thick, rounded corners=4pt,
        drop shadow={shadow xshift=0.12cm, shadow yshift=-0.12cm, opacity=0.12}
    },
    node_output/.style = {
        global_font, rectangle, draw=neutralgrey, fill=backgrey, thick, dashed, rounded corners=4pt
    },
    node_alert/.style = {
        global_font, rectangle, draw=alertred, fill=lightred, thick, rounded corners=4pt,
        drop shadow={shadow xshift=0.12cm, shadow yshift=-0.12cm, opacity=0.12}
    },
    % --- NODI SPECIALI PER INGEGNERIA INFORMATICA ---
    node_decision/.style = {
        global_font, diamond, aspect=2, draw=accentorange, fill=lightorange, thick,
        inner sep=3pt, drop shadow={shadow xshift=0.12cm, shadow yshift=-0.12cm, opacity=0.12}
    },
    node_database/.style = {
        global_font, cylinder, shape border rotate=90, aspect=0.25,
        draw=mainblue, fill=lightblue, thick, inner xsep=8pt, inner ysep=4pt,
        drop shadow={shadow xshift=0.12cm, shadow yshift=-0.12cm, opacity=0.12}
    },
    node_junction/.style = {
        global_font, circle, draw=neutralgrey, fill=white, thick, inner sep=3pt
    },
    % --- RIQUADRI DI RAGGRUPPAMENTO (FIT) ---
    group_box/.style = {
        draw=neutralgrey!50, fill=backgrey!40, dashed, thick, rounded corners=6pt,
        inner xsep=10pt, inner ysep=10pt
    },
    group_highlight/.style = {
        draw=paramgreen!60, fill=lightgreen!30, dotted, thick, rounded corners=6pt,
        inner xsep=10pt, inner ysep=10pt
    },
    % --- FRECCE E CONNETTORI ---
    arrow_source/.style = {-{Stealth[scale=1.1]}, draw=mainblue, very thick},
    arrow_target/.style = {-{Stealth[scale=1.1]}, draw=accentorange, very thick},
    arrow_param/.style = {-{Stealth[scale=1.1]}, draw=paramgreen, thick, dashed},
    arrow_flow/.style = {-{Stealth[scale=1.1]}, draw=neutralgrey, thick},
    arrow_alert/.style = {-{Stealth[scale=1.1]}, draw=alertred, thick},
    arrow_bidi/.style = {-{Stealth[scale=1.1]}, draw=neutralgrey, thick}
}

% Testo mostrato nel piè di pagina di ogni slide
\unimorefootertext{LT L-08 - A.A. 2025/2026}

\title{\small Analisi della Trasferibilità di Network Intrusion Detection Systems basati su Machine Learning al Dominio Satellitare}
\author{Alessandro Marchesini}

\begin{document}

% ------------------------------------------------------------------ %
% 1) TITOLO
% ------------------------------------------------------------------ %
\UnimoreTitlePage
  {Analisi della Trasferibilità di NIDS basati\\su Machine Learning al Dominio Satellitare}
  {Prof. Mirco Marchetti}
  {Ing. Dimitri Galli}
  {Alessandro Marchesini\\Matricola 187179}
  {Discussione della Prova Finale -- A.A. 2025/2026}
  {Dipartimento di Ingegneria ``Enzo Ferrari''}

% ------------------------------------------------------------------ %
% 2) INTRODUZIONE — CONTESTO
% ------------------------------------------------------------------ %
\section{Introduzione}

\begin{frame}{NIDS tra Dominio Terrestre e Satellitare}
  \begin{columns}[T]
    \begin{column}{0.5\textwidth}
      \centering
      \resizebox{!}{4.5cm}{%
        \begin{tikzpicture}
          \node[node_primary] (earth) {
            \textbf{Dominio Sorgente}\\
            rete terrestre
          };

          \node[node_process, below=0.7cm of earth] (nids) {
            \textbf{NIDS + ML}\\
            apprendimento\\
            supervisionato
          };

          \node[node_accent, below=0.7cm of nids] (space) {
            \textbf{Dominio Target}\\
            scenario terrestre-spaziale
          };

          \draw[arrow_source] (earth) -- node[left, font=\scriptsize] {training} (nids);
          \draw[arrow_target] (nids) -- node[left, font=\scriptsize] {inference} (space);

          \node[node_alert, right=1cm of nids] (shift) {
            \textbf{Domain Shift}\\
            distribuzioni di\\
            traffico differenti
          };

          \draw[arrow_alert] (earth.east) to[bend left=20] (shift.north);
          \draw[arrow_alert] (shift.south) to[bend left=20] (space.east);
        \end{tikzpicture}
      }
    \end{column}

    \begin{column}{0.3\textwidth}
      \begin{block}{Limite}
        \footnotesize Le prestazioni \textbf{in-domain} non sono sufficienti per caratterizzare la robustezza di un NIDS in presenza di \textbf{domain shift}.
      \end{block}
    \end{column}
  \end{columns}
\end{frame}

% ------------------------------------------------------------------ %
% 3) PROBLEM STATEMENT — RQ
% ------------------------------------------------------------------ %
\section{Problem Statement}

\begin{frame}{Problema di Ricerca e Domande Sperimentali}
  \centering
  \begin{columns}[T]
    \begin{column}{0.3\textwidth}
      \begin{block}{Obiettivo}
        \footnotesize Valutare la \textbf{generalizzazione cross-domain} di modelli addestrati sul dominio sorgente terrestre verso distribuzioni target terrestre-spaziali.      
      \end{block}
    \end{column}

    \begin{column}{0.6\textwidth}
      \begin{tikzpicture}
          \draw[draw=unimoreMaroon, dashed, thick]
              plot[smooth] coordinates {
                  (0,0)
                  (0,-1)
                  (0,-2)
                  (0,-3)
              };

          % Punti rossi
          \fill[unimoreMaroon] (0,0) circle (3pt);
          \fill[unimoreMaroon] (0,-1) circle (3pt);
          \fill[unimoreMaroon] (0,-2) circle (3pt);
          \fill[unimoreMaroon] (0,-3) circle (3pt);

          % Etichette
          \node[align=left, text=unimoreMaroon, anchor=west] at (0.2,0)
              {\textbf{RQ1} \footnotesize Trasferibilità diretta e Domain Shift};

          \node[align=left, text=unimoreMaroon, anchor=west] at (0.2,-1)
              {\textbf{RQ2} \footnotesize Adattamento mediante Conoscenza Target Sparsa};

          \node[align=left, text=unimoreMaroon, anchor=west] at (0.2,-2)
              {\textbf{RQ3} \footnotesize Granularità, Specializzazione e Generalizzazione};

          \node[align=left, text=unimoreMaroon, anchor=west] at (0.2,-3)
              {\textbf{RQ4} \footnotesize Architettura, Stabilità e Struttura Decisionale};

      \end{tikzpicture}
    \end{column}
  \end{columns}
\end{frame}

% ------------------------------------------------------------------ %
% 4) CONTRIBUTI / FRAMEWORK
% ------------------------------------------------------------------ %
\section{Contributi}

\begin{frame}{Framework di Valutazione Cross-Domain}
  \centering
  \resizebox{!}{3cm}{
    \begin{tikzpicture}
      \node[node_output] (analysis) {
        \textbf{Diagnostica e interpretazione}\\
        PCA / KDE \;\; SHAP \;\; TPR--TNR--F1--AP \;\; test $t$ di Welch
      };

      \node[node_process, above=1cm of analysis, xshift=0.6cm] (train) {
        \textbf{Training}\\
        \textbf{Controllato}
      };

      \node[node_primary, left=0.7cm of train] (repr) {
        \textbf{Rappresentazione}\\
        \textbf{Comune}
      };

      \node[node_accent, right=0.7cm of train] (eval) {
        \textbf{Cross}\\
        \textbf{Evaluation}
      };

      \draw[arrow_flow] (repr) -- (train);
      \draw[arrow_flow] (train) -- (eval);

      

      \draw[arrow_flow] (repr.south) -- ($(analysis.north) + (-3,0)$);
      \draw[arrow_param] (train.south) -- ($(analysis.north) + (0.6,0)$);
      \draw[arrow_flow] (eval.south) -- ($(analysis.north) + (3.56,0)$);
    \end{tikzpicture}
  }

  \vspace{0.30cm}

  \footnotesize Il contributo consiste in un \textbf{protocollo sistematico per studiare la trasferibilità} di NIDS supervisionati tra domini eterogenei.
\end{frame}

% ------------------------------------------------------------------ %
% 5) DATASET E MAPPING
% ------------------------------------------------------------------ %
\section{Metodologia}

\begin{frame}{Preprocessing \& Datasets}
  \begin{columns}[T]
    \begin{column}{0.5\textwidth}
      \centering
      \resizebox{!}{4.5cm}{
        \begin{tikzpicture}
          \node [node_transform] (align) {
              \textbf{Operatore di Mappatura e Pulizia}
          };
      
          % --- FASE 1: INGESTIONE E DATASET GREZZI ---
          \node [node_database, above=1cm of align, xshift=-2cm, minimum width=3.2cm, minimum height=2cm] (unsw) {
              \textbf{UNSW-NB15}\\
              \footnotesize $\mathcal{D}_S$
          };
      
          \node [node_database, draw=accentorange, fill=lightorange, above=1cm of align, xshift=2cm, minimum width=3.2cm, minimum height=2cm] (stin) {
              \textbf{STIN}\\
              \footnotesize $\mathcal{D}_{T \in \{TER20-SAT20\}}$
          };
      
          \node [node_transform, below=1cm of align] (sampling) {
              \textbf{Campionamento Stratificato e Partizionamento}\\
              $ben/mal = 10 \; (10:1)$ \;\; $tr/te = 80/20$\\
              5 Benchmark Aggregati + 11 Benchmark Biclasse
          };

          \draw [arrow_source] (unsw.south) -- (align.north);
          \draw [arrow_target] (stin.south) -- (align.north);
          \draw [arrow_flow] (align) -- (sampling);
        \end{tikzpicture}
      }
    \end{column}

    \begin{column}{0.5\textwidth}
      \centering
      \resizebox{!}{4.5cm}{
        \begin{tikzpicture}
          \node[node_primary, minimum width=5cm] (sourceModel) {
            \textbf{Source-only}\\
            \\
            Ben: $100\%\mathcal{D}_S$\\
            Mal: $100\%\mathcal{D}_S$
          };

          \node[node_process, below=0.3cm of sourceModel, minimum width=5cm] (injModel) {
            \textbf{INJECTION}\\
            \\
            Ben: $100\%\mathcal{D}_S$\\
            Mal: $\approx 99.7\%\mathcal{D}_S + 0.3\%\mathcal{D}_T$
          };

          \node[node_accent, below=0.3cm of injModel, minimum width=5cm] (hybridModel) {
            \textbf{Hybrid}\\
            \\
            Ben: $100\%\mathcal{D}_S$\\
            Mal: $100\%\mathcal{D}_T$
          };

          \node [right, font=\footnotesize, text=neutralgrey, align=left] at (sourceModel.east) {
            Conoscenza target\\
            \textbf{assente}
          };
          \node [right, font=\footnotesize, text=paramgreen, align=left] at (injModel.east) {
            Conoscenza target\\
            \textbf{sparsa}
          };
          \node [right, font=\footnotesize, text=accentorange, align=left] at (hybridModel.east) {
            Conoscenza target\\
            \textbf{dominante}
          };
        \end{tikzpicture}
      }
    \end{column}
  \end{columns}
\end{frame}

% ------------------------------------------------------------------ %
% 6) TRE PARADIGMI
% ------------------------------------------------------------------ %
\begin{frame}{Models \& Validation}
  \begin{columns}
    \begin{column}{0.5\textwidth}
      \centering
      \resizebox{\textwidth}{!}{
        \begin{tikzpicture}[taxonomy_dot/.style={circle, fill=unimoreMaroon, inner sep=2.8pt}]
            % ---------------------------------------------------------------------
            % NODI
            % ---------------------------------------------------------------------
            \node[taxonomy_dot, label={[anchor=west]right:\textbf{TREE-BASED METHODS}}] (tree) at (0,0) {};

            \node[taxonomy_dot, label={[anchor=west]right:\textbf{SINGLE TREE}}] (single) at (-2.8,-2.8) {};
            \node[taxonomy_dot, label={[anchor=west]right:\textbf{ENSEMBLE}}] (ensemble) at (2.8,-2.8) {};

            \node[taxonomy_dot, label={[anchor=west]right:\textbf{BAGGING}}] (bagging) at (1.4,-4.4) {};
            \node[taxonomy_dot, label={[anchor=west]right:\textbf{BOOSTING}}] (boosting) at (4.2,-4.4) {};

            \node[taxonomy_dot, label={[anchor=west]right:\textbf{DT}}] (dt) at (-2.8,-6.0) {};
            \node[taxonomy_dot, label={[anchor=west]right:\textbf{RF}}] (rf) at (1.4,-6.0) {};
            \node[taxonomy_dot, label={[anchor=west]right:\textbf{HGB}}] (hgb) at (4.2,-6.0) {};

            % ---------------------------------------------------------------------
            % CONNETTORI
            % ---------------------------------------------------------------------
            % TREE-BASED METHODS -> SINGLE TREE / ENSEMBLE
            \draw[thick] (tree) -- (0,-2.0);
            \draw[thick] (-2.8,-2.0) -- (2.8,-2.0);
            \draw[thick] (-2.8,-2.0) -- (single);
            \draw[thick] (2.8,-2.0) -- (ensemble);

            % SINGLE TREE -> DT
            \draw[thick] (single) -- (dt);

            % ENSEMBLE -> BAGGING / BOOSTING
            \draw[thick] (ensemble) -- (2.8,-3.6);
            \draw[thick] (1.4,-3.6) -- (4.2,-3.6);
            \draw[thick] (1.4,-3.6) -- (bagging);
            \draw[thick] (4.2,-3.6) -- (boosting);

            % BAGGING -> RF
            \draw[thick] (bagging) -- (rf);

            % BOOSTING -> HGB
            \draw[thick] (boosting) -- (hgb);
        \end{tikzpicture}
      }
    \end{column}

    \begin{column}{0.5\textwidth}
      \centering
      \resizebox{!}{5.5cm}{
        \begin{tikzpicture}[]

            % Livello 1: Dati Iniziali
            \node [node_primary] (data) {Traffico Preprocessato};

            % Livello 3: Nodi Multi-seed e Addestramenti Paralleli
            \node [node_accent, below=0.7cm of data] (seed2) {
                Seed\\
                Stocastico $s_2$
            };
            \node [node_accent, left=1cm of seed2] (seed1) {
                Seed\\
                Stocastico $s_1$
            };
            \node [right=0.15cm of seed2] (dots1) {\dots};
            \node [node_accent, right=1cm of seed2] (seedK) {
                Seed\\
                Stocastico $s_K$
            };

            \node [node_process, below=0.7cm of seed1] (train1) {
                Addestramento\\
                Modello $M$
            };
            \node [node_process, below=0.7cm of seed2] (train2) {
                Addestramento\\
                Modello $M$
            };
            \node [right=0cm of train2] (dots2) {\dots};
            \node [node_process, below=0.7cm of seedK] (trainK) {
                Addestramento\\
                Modello $M$
            };

            % Frecce dal blocco Benchmark ai percorsi paralleli
            \draw [arrow_flow] ($(data.south) + (0, -0.35)$) -| (seed1);
            \draw [arrow_flow] (data.south) -- (seed2);
            \draw [arrow_flow] ($(data.south) + (0, -0.35)$) -| (seedK);

            \draw [arrow_flow] (seed1) -- (train1);
            \draw [arrow_flow] (seed2) -- (train2);
            \draw [arrow_flow] (seedK) -- (trainK);

            % Livello 4: Valutazione
            \node [node_output, below=0.7cm of train1] (eval1) {
                Matrice $H_1$
            };
            \node [node_output, below=0.7cm of train2] (eval2) {
                Matrice $H_2$
            };
            \node [right=0.3cm of eval2] (dots3) {\dots};            
            \node [node_output, below=0.7cm of trainK] (evalK) {
                Matrice $H_K$
            };

            \draw [arrow_flow] (train1) -- (eval1);
            \draw [arrow_flow] (train2) -- (eval2);
            \draw [arrow_flow] (trainK) -- (evalK);

            % Livello 5: Aggregazione Statistica
            \node [node_process, below=0.7cm of eval2] (agg) {
                Calcolo Stime Campionarie
            };
            
            \draw [arrow_flow] (eval1.south) |- ++(0, -0.35) -| (agg.north);
            \draw [arrow_flow] (eval2) -- (agg);
            \draw [arrow_flow] (evalK.south) |- ++(0,-0.35) -| (agg.north);

            % Livello 6: Test
            \node [node_output, below=0.7cm of agg] (welch) {
                Test $t$ di Welch ($p$-value)
            };
            
            \draw [arrow_flow] (agg) -- (welch);
        \end{tikzpicture}
      }
    \end{column}
  \end{columns}
\end{frame}

% ------------------------------------------------------------------ %
% 7) IMPLEMENTAZIONE / PROTOCOLLO
% ------------------------------------------------------------------ %
\section{Implementazione e Risultati}

\begin{frame}{Implementazione e Protocollo Sperimentale}
  \centering
  \begin{tikzpicture}[node distance=0.55cm]
    \node[node_database, text width=1.75cm, minimum height=1.05cm] (data) {
      \textbf{Dataset}\\
      NB15 / STIN
    };

    \node[node_transform, right=0.45cm of data, text width=2.05cm, minimum height=1.05cm] (prep) {
      \textbf{Preprocessing}\\
      mapping a 16 feature
    };

    \node[node_output, right=0.45cm of prep, text width=2.05cm, minimum height=1.05cm] (bench) {
      \textbf{Benchmark}\\
      5 aggregati + 11 biclasse
    };

    \node[node_process, right=0.45cm of bench, text width=1.85cm, minimum height=1.05cm] (trainp) {
      \textbf{Training}\\
      DT / HGB / RF
    };

    \node[node_accent, right=0.45cm of trainp, text width=1.85cm, minimum height=1.05cm] (evalp) {
      \textbf{Evaluation}\\
      16 test set
    };

    \draw[arrow_flow] (data) -- (prep);
    \draw[arrow_flow] (prep) -- (bench);
    \draw[arrow_param] (bench) -- (trainp);
    \draw[arrow_target] (trainp) -- (evalp);
  \end{tikzpicture}

  \vspace{0.45cm}

  \begin{columns}[T]
    \begin{column}{0.32\textwidth}
      \begin{block}{Riproducibilità}
        \centering
        $K=10$ seed\\[1mm]
        stesso protocollo e stessi test set
      \end{block}
    \end{column}

    \begin{column}{0.32\textwidth}
      \begin{block}{Valutazione}
        \centering
        TPR / TNR / Precision / F1 / AP\\[1mm]
        PCA / KDE / SHAP
      \end{block}
    \end{column}

    \begin{column}{0.32\textwidth}
      \begin{block}{Analisi statistica}
        \centering
        macro-F1 sui 16 benchmark\\[1mm]
        test $t$ di Welch esplorativo
      \end{block}
    \end{column}
  \end{columns}
\end{frame}

% ------------------------------------------------------------------ %
% 8) RISULTATO PRINCIPALE — SOURCE VS INJECTION
% ------------------------------------------------------------------ %
\begin{frame}{Risultato Principale: Effetto della Conoscenza Target Sparsa}
  \begin{columns}[T]
    \begin{column}{0.67\textwidth}
      \centering
      \textbf{TPR media sui nove benchmark target}

      \vspace{0.25cm}
      \begin{tikzpicture}[x=6.0cm,y=0.60cm]
        % Asse e griglia
        \draw[arrow_flow] (0,0) -- (1.06,0) node[right, font=\scriptsize] {TPR};
        \foreach \x/\lab in {0/0,0.25/0.25,0.5/0.50,0.75/0.75,1/1.00} {
          \draw[neutralgrey!20] (\x,0.12) -- (\x,6.35);
          \node[below, font=\scriptsize] at (\x,0) {\lab};
        }

        % DT
        \node[left, font=\small\bfseries] at (0,5.65) {DT};
        \fill[mainblue!75] (0,5.35) rectangle (0.3234,5.65);
        \fill[paramgreen!80] (0,4.90) rectangle (0.96612,5.20);
        \node[right, font=\scriptsize] at (0.3234,5.50) {0.3234};
        \node[right, font=\scriptsize\bfseries] at (0.96612,5.05) {0.9661};

        % HGB
        \node[left, font=\small\bfseries] at (0,3.85) {HGB};
        \fill[mainblue!75] (0,3.55) rectangle (0.0607,3.85);
        \fill[paramgreen!80] (0,3.10) rectangle (0.82970,3.40);
        \node[right, font=\scriptsize] at (0.0607,3.70) {0.0607};
        \node[right, font=\scriptsize\bfseries] at (0.82970,3.25) {0.8297};

        % RF
        \node[left, font=\small\bfseries] at (0,2.05) {RF};
        \fill[mainblue!75] (0,1.75) rectangle (0.3035,2.05);
        \fill[paramgreen!80] (0,1.30) rectangle (0.98821,1.60);
        \node[right, font=\scriptsize] at (0.3035,1.90) {0.3035};
        \node[right, font=\scriptsize\bfseries] at (0.98821,1.45) {0.9882};

        % Legenda
        \fill[mainblue!75] (0.03,0.45) rectangle (0.075,0.68);
        \node[anchor=west, font=\scriptsize] at (0.085,0.565) {Source-only};
        \fill[paramgreen!80] (0.36,0.45) rectangle (0.405,0.68);
        \node[anchor=west, font=\scriptsize] at (0.415,0.565) {INJECTION};
      \end{tikzpicture}
    \end{column}

    \begin{column}{0.33\textwidth}
      \begin{exampleblock}{Solo 29 campioni target}
        Incremento TPR target:
        \begin{itemize}
          \item DT: $+0.6427$
          \item HGB: $+0.7690$
          \item RF: $+0.6847$
        \end{itemize}
      \end{exampleblock}

      \begin{block}{Prestazioni sorgente}
        La TPR media sui benchmark NB15 rimane sostanzialmente invariata:
        \[
        |\Delta \mathrm{TPR}| \approx 10^{-3}.
        \]
      \end{block}
    \end{column}
  \end{columns}
\end{frame}

% ------------------------------------------------------------------ %
% 9) RISULTATI GLOBALI — TRADE-OFF
% ------------------------------------------------------------------ %
\begin{frame}{Risultati Globali: Generalizzazione vs Specializzazione}
  \begin{columns}[T]
    \begin{column}{0.60\textwidth}
      \centering
      \textbf{TPR media sull'intera cross-evaluation (16 benchmark)}

      \vspace{0.25cm}
      \small
      \renewcommand{\arraystretch}{1.35}
      \begin{tabular}{lccc}
        \toprule
        \textbf{Paradigma} & \textbf{DT} & \textbf{HGB} & \textbf{RF} \\
        \midrule
        Source-only & 0.5288 & 0.3694 & 0.5161 \\
        \textbf{INJECTION} & \textbf{0.8899} & \textbf{0.8025} & \textbf{0.9007} \\
        Hybrid aggregato & 0.5626 & 0.5525--0.5627 & 0.5626 \\
        \bottomrule
      \end{tabular}

      \vspace{0.38cm}
      \begin{tikzpicture}[node distance=0.35cm]
        \node[node_primary, text width=2.55cm, minimum height=1.35cm] (s) {
          \textbf{Source-only}\\
          buono su NB15\\
          debole sul target
        };

        \node[node_process, right=0.35cm of s, text width=2.55cm, minimum height=1.35cm] (i) {
          \textbf{INJECTION}\\
          risposta più uniforme\\
          source + target
        };

        \node[node_accent, right=0.35cm of i, text width=2.55cm, minimum height=1.35cm] (h) {
          \textbf{Hybrid}\\
          TPR target $\approx 1$\\
          TPR NB15 $\approx 0$
        };
      \end{tikzpicture}
    \end{column}

    \begin{column}{0.40\textwidth}
      \begin{alertblock}{Specializzazione ibrida}
        \small
        TPR target: circa $0.9818$--$0.9999$.\\[1mm]
        TPR sugli attacchi NB15: quasi nulla.
      \end{alertblock}

      \begin{exampleblock}{Architetture nella baseline}
        \small
        \textbf{RF}: miglior profilo predittivo misurato
        \[
        \mathrm{TPR}_{target}=0.98821,
        \quad
        \mathrm{F1}=0.890221
        \]
        \textbf{HGB}: TNR più elevata, $0.99658$.
      \end{exampleblock}
    \end{column}
  \end{columns}
\end{frame}

% ------------------------------------------------------------------ %
% 10) CONCLUSIONI
% ------------------------------------------------------------------ %
\section{Conclusioni}

\begin{frame}{Conclusioni}
  \centering
  \begin{tikzpicture}[node distance=0.55cm]
    \node[node_primary, text width=3.25cm, minimum height=2.15cm] (c1) {
      \textbf{1. Trasferibilità diretta}\\[1mm]
      Limitata e dipendente dalla distribuzione target.
    };

    \node[node_process, right=0.65cm of c1, text width=3.25cm, minimum height=2.15cm] (c2) {
      \textbf{2. Conoscenza target sparsa}\\[1mm]
      INJECTION è associata a una forte crescita della sensibilità target preservando il sorgente.
    };

    \node[node_accent, right=0.65cm of c2, text width=3.25cm, minimum height=2.15cm] (c3) {
      \textbf{3. Specializzazione}\\[1mm]
      Massimizzare il target non equivale a massimizzare la robustezza cross-domain.
    };
  \end{tikzpicture}

  \vspace{0.45cm}

  \begin{block}{Limiti da mantenere espliciti}
    \centering
    Benign-domain shift non valutato
    \quad $\bullet$ \quad
    target basato su STIN
    \quad $\bullet$ \quad
    29 campioni non sono una quantità minima/ottimale
    \quad $\bullet$ \quad
    costi di deployment non misurati
  \end{block}

  \vspace{0.25cm}
  \Large
  \textbf{La prestazione intra-domain non è sufficiente per valutare\\
  la robustezza di un NIDS destinato a scenari eterogenei.}
\end{frame}

\end{document}
